\begin{examplebox}
\textbf{\textcolor{CExample}{例 1（基础）}}

\textbf{\textcolor{CExample}{题目：}} 比较 $\sin 1$ 与 $\frac{2}{\pi}$ 的大小。

\textbf{\textcolor{CExample}{解题步骤：}}
\begin{description}[style=nextline, leftmargin=2.8em, labelsep=0.5em, itemsep=0.45em]
    \item[\textbf{第一步（确认范围）：}] 判断数值 $1$ 是否在结论适用区间内。
    \par\textbf{理由：} 结论要求 $x \in [0, \frac{\pi}{2}]$。因为 $0 < 1 < \frac{\pi}{2} \approx 1.57$，所以 $x=1$ 在区间内。

    \item[\textbf{第二步（直接应用）：}] 将 $x=1$ 代入结论不等式。
    \[
    \sin 1 \geq \frac{2}{\pi} \times 1 = \frac{2}{\pi}
    \]
    \par\textbf{理由：} 根据命题 $\sin x \geq \frac{2}{\pi}x$，令 $x=1$ 即得。

    \item[\textbf{第三步（判断大小）：}] 由不等式直接得到比较结果。
    \par\textbf{理由：} 不等式 $\sin 1 \geq \frac{2}{\pi}$ 已经表明 $\sin 1$ 不小于 $\frac{2}{\pi}$。由于 $1$ 不是区间端点，等号不成立，故严格大于。
\end{description}

\textbf{\textcolor{CExample}{关键结论：}} \textbf{$\sin 1 > \frac{2}{\pi}$。}

\medskip
\textbf{\textcolor{CExample}{例 2（稍变形）}}

\textbf{\textcolor{CExample}{题目：}} 已知 $\alpha \in (0, \frac{\pi}{2})$，求证：$\frac{\sin \alpha}{\alpha} > \frac{2}{\pi}$。

\textbf{\textcolor{CExample}{解题步骤：}}
\begin{description}[style=nextline, leftmargin=2.8em, labelsep=0.5em, itemsep=0.45em]
    \item[\textbf{第一步（分离目标）：}] 将待证不等式两边同乘以正数 $\alpha$，进行变形。
    \[
    \frac{\sin \alpha}{\alpha} > \frac{2}{\pi} \quad \Leftrightarrow \quad \sin \alpha > \frac{2}{\pi} \alpha
    \]
    \par\textbf{理由：} 因为 $\alpha \in (0, \frac{\pi}{2})$，所以 $\alpha > 0$，不等号方向不变。

    \item[\textbf{第二步（匹配结论）：}] 观察变形后的不等式，与已知结论形式一致。
    \par\textbf{理由：} 变形后即需证 $\sin \alpha > \frac{2}{\pi} \alpha$，其中 $\alpha \in (0, \frac{\pi}{2})$。

    \item[\textbf{第三步（应用结论）：}] 直接应用已知命题，并说明等号不成立。
    \par\textbf{理由：} 由命题知，当 $\alpha \in [0, \frac{\pi}{2}]$ 时，有 $\sin \alpha \geq \frac{2}{\pi} \alpha$。
    \par 由于 $\alpha \in (0, \frac{\pi}{2})$，不是区间端点 $0$ 或 $\frac{\pi}{2}$，因此等号不成立，故有严格不等式 $\sin \alpha > \frac{2}{\pi} \alpha$ 成立。
\end{description}

\textbf{\textcolor{CExample}{关键结论：}} \textbf{因此，原不等式 $\frac{\sin \alpha}{\alpha} > \frac{2}{\pi}$ 得证。}
\end{examplebox}
